\documentclass[11pt]{article}
\usepackage{amsmath, amssymb}
\usepackage{booktabs}
\usepackage{tikz-cd}
\usepackage[margin=1in]{geometry}

\title{The Geometry of Data Is the Program}
\author{PhaseShift Studio}
\date{}

\begin{document}
\maketitle

\begin{abstract}
A function's type is usually read as a constraint on its values. In a
data-oriented machine it is something stronger: a description of the
\emph{geometry} of the data being walked---and therefore of the control flow
itself. Where the data branches, the code parallelizes; where it merges, the
code synchronizes. The signature carries both.
\end{abstract}

\section{The coefficient is the geometry}

Multiplicity lives in the type as a \emph{coefficient}: a $(min, max)$ pair on
every type. Four shapes do the work:

\begin{center}
\begin{tabular}{@{}ccc@{}}
\toprule
coefficient & meaning & form \\
\midrule
$\{0\}$ & nothing      & seed \\
$\{?\}$ & zero or one  & maybe / filter \\
$\{1\}$ & exactly one  & map \\
$\{*\}$ & any number   & gather / scatter \\
\bottomrule
\end{tabular}
\end{center}

An instruction is a morphism $f : A^{c_{\mathrm{dom}}} \to B^{c_{\mathrm{rng}}}$,
and its \emph{form} is read off the coefficient pair:

\begin{center}
\begin{tabular}{@{}ccc@{}}
\toprule
form    & signature        & what it is \\
\midrule
map     & $\{1\} \to \{1\}$ & one in, one out --- pointwise, freely parallel \\
scatter & $\{1\} \to \{*\}$ & one in, many out --- a \emph{branch} \\
reduce  & $\{*\} \to \{1\}$ & many in, one out --- a \emph{merge} \\
seed    & $\{0\} \to A$     & from nothing \\
\bottomrule
\end{tabular}
\end{center}

\section{The tree walk is the schedule}

Data lives at addresses, and addresses are paths in a tree. Descending a node
with $N$ children \emph{splits} the traversal---the multiplicity grows to $N$.
Ascending \emph{merges} it, and the coefficient is the bulk that carries the
count back up. The two directions are the two fundamental operations:

\begin{center}
\begin{tikzcd}
	& x \arrow[dl] \arrow[d] \arrow[dr] & \\
	a & b & c
\end{tikzcd}
\qquad
\begin{tikzcd}
	a \arrow[dr] & b \arrow[d] & c \arrow[dl] \\
	& y &
\end{tikzcd}

\small scatter $\{1\} \to \{*\}$ (left) and gather $\{*\} \to \{1\}$ (right)
\end{center}

Descending \emph{is} the fan-out; ascending \emph{is} the fan-in. The data
structure hands the compiler its own shape: the branching factor becomes the
coefficient, and the coefficient is the schedule.

\section{Category theory, in one line}

Scatter is a \emph{cocartesian fan-out} (one wire, many); gather is the
\emph{limit} that reunifies them. The coefficient is a \emph{measure on the
tree}---the multiplicity of the fiber over a point---and an instruction's form
is how that measure changes across the morphism.

\begin{quote}
The compiler doesn't infer parallelism from control flow; it reads it off the
data's geometry.
\end{quote}

Branch $\Rightarrow$ parallelize. Merge $\Rightarrow$ barrier. Coefficient
$\Rightarrow$ signature. No annotations, no static control-flow analysis---the
data carries the plan.

\end{document}
